% Line Atlas chapter draft.
% Suggested dependencies: amsmath, amssymb, siunitx, booktabs, tabularx,
% tikz, pgfplots, xcolor, hyperref.

\chapter{Stellar Feedback, Superbubbles, and Galactic Outflows}
\label{chap:feedback-outflows}
\chaptermark{Stellar Feedback and Galactic Outflows}

Massive stars reshape their surroundings before, during, and after their deaths.
Photoionization heats gas; radiation transfers momentum; winds and jets inflate hot
cavities; supernovae generate shocks, metals, cosmic rays, and turbulence.  These
channels overlap in space and time.  Their collective action can disperse molecular
clouds, maintain a multiphase disk, launch fountains and winds, and seed the
circumgalactic medium.

``Feedback'' is not a synonym for broad lines.  It is a causal claim that a source has
transferred energy, momentum, mass, or newly synthesized material to gas.  Establishing
that claim requires a budget, geometry, and timescale.  Reviews spanning local and
galactic scales include \citet{McKeeOstriker2007}, \citet{NaabOstriker2017},
\citet{VeilleuxCecilBlandHawthorn2005}, and
\citet{VeilleuxMaiolinoBolattoAalto2020}.

\section{Bookkeeping: rates, integrated yields, and coupling}

For a feedback channel with injection rate $\dot E_{\rm inj}$, only a fraction
$\epsilon_E$ may appear as kinetic plus thermal energy of the observed phase:
\begin{equation}
 \dot E_{\rm phase}=\epsilon_E\dot E_{\rm inj}
 -\dot E_{\rm rad}-\dot E_{\rm transfer}-\dot E_{\rm work}.
 \label{eq:feedback-energy-budget}
\end{equation}
Similarly, momentum and mass move between phases.  A low energy measured in cool gas
does not demonstrate weak feedback if most energy resides in a faint hot wind.
Conversely, comparing a cool phase's momentum with a source's instantaneous rate can
mislead if the flow accumulated over a long episode.

The useful quantities have distinct dimensions:
\begin{align}
 \dot M_{\rm out} &\quad [M\,T^{-1}],\\
 \dot p_{\rm out}=\dot M_{\rm out}v &\quad [M\,L\,T^{-2}],\\
 \dot E_{\rm kin}=\tfrac12\dot M_{\rm out}v^2 &\quad [M\,L^2\,T^{-3}].
 \label{eq:outflow-fluxes}
\end{align}
Papers sometimes call $p=Mv$ a ``momentum rate'' or call $E/t$ ``energy.''  Dimensional
labels and equations prevent such ambiguities.

Define mass loading relative to star formation as
\begin{equation}
 \eta_M=\frac{\dot M_{\rm out}}{\dot M_\star},
 \label{eq:mass-loading}
\end{equation}
with analogous metal or energy loading factors.  Numerator and denominator should
refer to compatible times and spatial boundaries.  An outflow at 10~kpc records a
past launch episode, whereas H$\alpha$ may weight the most recent few Myr.

\section{Photoionization heating and expansion}

Each hydrogen photoionization leaves the excess photon energy above 13.6~eV as
photoelectron energy, which heats the gas.  Thermal balance near $10^4$~K raises the
pressure relative to cold molecular surroundings, driving an ionization front and
shock.  The Str\"omgren and ideal expansion solutions in
Chapter~\ref{chap:star-formation} provide baseline scales.

Photoionization can both suppress and promote local collapse.  It erodes and
accelerates cloud surfaces, lowers the mass available to form stars, and opens
low-density escape channels.  Its shock can also compress pre-existing dense
structures.  Demonstrating ``triggered'' star formation requires showing that stars
would not have formed on the same timescale without the feedback, a counterfactual
rarely established by spatial coincidence alone.

Optical line ratios identify ionization fronts and partially ionized zones.  Enhanced
[S\,II]/H$\alpha$ or [O\,I]/H$\alpha$ can arise in shocks, diffuse ionized gas, or hard
radiation.  Kinematics and multiple diagnostic planes are required.  Free--free and
recombination emission trace ionized gas even when optical emission is obscured.

\section{Radiation pressure}

Absorption of luminosity $L$ supplies a direct momentum rate $L/c$.  A schematic
radiative force is
\begin{equation}
 \dot p_{\rm rad}\simeq
 \left(1-e^{-\tau_{\rm UV}}+f_{\rm trap}\tau_{\rm IR}\right)\frac{L}{c}.
 \label{eq:radiation-momentum}
\end{equation}
The first term accounts for direct absorption; the second is an idealized infrared
multiple-interaction contribution.  The factor $f_{\rm trap}$ recognizes that
porosity, radiation--matter instabilities, and geometry can reduce coupling below a
homogeneous optical-depth estimate.  It should be constrained, not assumed to be one.

Resonance-line photons also transfer momentum through repeated scattering, but dust
destruction, frequency diffusion, and escape through velocity or spatial channels
limit the number of effective interactions.  The emergent line profile is a probe of
the same transfer process and cannot be interpreted as an independent intrinsic
emission spectrum.

Comparing an observed momentum flux with $L/c$ is useful, but a ratio above unity does
not uniquely prove radiation trapping.  Energy-conserving hot gas, time variability,
geometry, radius estimates, and phase conversion can all produce an apparent
``momentum boost.''

\section{Stellar winds and wind-blown bubbles}

A stellar wind with mass-loss rate $\dot M_w$ and terminal speed $v_w$ carries
\begin{equation}
 \dot E_w=\frac12\dot M_wv_w^2,
 \qquad \dot p_w=\dot M_wv_w.
 \label{eq:stellar-wind-power}
\end{equation}
The freely expanding wind terminates at a reverse shock, creating a hot interior that
pushes a swept-up shell.  For constant mechanical luminosity $L_w$ in a uniform
ambient density $\rho_0$, the ideal energy-conserving similarity solution scales as
\begin{equation}
 R_b(t)=C\left(\frac{L_w}{\rho_0}\right)^{1/5}t^{3/5},
 \qquad v_b=\frac{3R_b}{5t},
 \label{eq:wind-bubble-scaling}
\end{equation}
where $C$ depends on the exact solution convention \citep{WeaverEtAl1977}.  Radiative
cooling, conduction, turbulent mixing, leakage, density gradients, and episodic input
alter both normalization and exponent.

Bubble interiors can be faint because emissivity scales roughly as density squared.
Soft X-ray luminosity often probes dense interfaces or mixed material rather than the
volume-filling wind.  Optical shells trace cooled swept-up gas, and expansion may be
seen as split profiles or coherent velocity ellipses.  A ring on the sky without
kinematic expansion is not sufficient evidence for a bubble.

\subsection{Worked example: distinguishing energy- and momentum-driven scalings}

An energy-driven bubble in uniform density has $R\propto t^{3/5}$, while a shell
driven by constant momentum injection scales approximately as $R\propto t^{1/2}$.
Suppose a bubble has $R=100$~pc and $v=30\,\mathrm{km\,s^{-1}}$.  The similarity ages
are
\begin{equation}
 t_E=\frac{3R}{5v}\simeq2.0~\mathrm{Myr},\qquad
 t_p=\frac{R}{2v}\simeq1.6~\mathrm{Myr}.
\end{equation}
The difference is smaller than typical uncertainties from projection and acceleration
through a density gradient.  Radius and speed alone rarely establish the driving
regime; hot-gas pressure, shell mass, energy losses, and source history add leverage.

\section{Supernova remnants and terminal momentum}

A supernova injects ejecta mass, metals, and energy of order $10^{51}$~erg for a
canonical energetic event.  Early free expansion gives way to a Sedov--Taylor stage
when swept-up mass dominates.  For explosion energy $E$ in uniform density $\rho_0$,
\begin{equation}
 R_{\rm ST}(t)=\xi\left(\frac{Et^2}{\rho_0}\right)^{1/5},
 \qquad v_{\rm ST}=\frac{2R_{\rm ST}}{5t}.
 \label{eq:sedov-taylor}
\end{equation}
When shocked gas cools efficiently, a dense radiative shell forms.  Energy declines,
but radial momentum continues to grow modestly before approaching a terminal value.
Numerical calculations for representative conditions find a weak density dependence
and terminal momentum of order a few
$10^5\,M_\odot\,\mathrm{km\,s^{-1}}$ per event, with environment, clustering, and
microphysics affecting the result \citep{KimOstriker2015}.

An unresolved simulation that deposits thermal energy into too much mass can radiate
it immediately and miss the physical momentum-generating stage.  Feedback
implementations therefore require adequate resolution or a calibrated terminal
momentum prescription.  Injecting both full energy and terminal momentum without
accounting for the resolved stage can double count.

Observed remnants span nonradiative X-ray and Balmer shocks, radiative optical shells,
radio synchrotron structures, and molecular interactions.  Shock speeds inferred from
line ratios or X-ray temperature should be checked against proper-motion and line-width
measurements.  Electron and ion temperatures may not equilibrate immediately behind a
fast collisionless shock.

\section{Clustering, superbubbles, and breakout}

Massive stars form in groups, so winds and supernovae overlap.  Explosions inside an
existing low-density cavity lose less energy than isolated explosions in dense gas,
while shell impacts and mixing can enhance cooling.  The interval between events,
cluster compactness, ambient stratification, and pre-supernova clearing determine
whether energy accumulates coherently.

In a vertically stratified disk, a superbubble may accelerate as it reaches lower
density, fragment, and vent hot gas into the halo.  Breakout connects local feedback
to a galactic fountain or wind.  The dense shell may remain near the disk while hot,
metal-rich gas escapes through channels, so mass, metal, and energy loading need not
track each other.

\begin{figure}[tbp]
\centering
\begin{tikzpicture}[
 node distance=9mm and 8mm,
 every node/.style={font=\small},
 source/.style={draw,rounded corners,fill=yellow!18,align=center,minimum width=25mm},
 phase/.style={draw,rounded corners,fill=blue!9,align=center,minimum width=27mm},
 arrow/.style={->,thick}
]
 \node[source] (stars) {massive stars\\and clusters};
 \node[source,below left=of stars] (rad) {ionizing and\\nonionizing radiation};
 \node[source,below=of stars] (winds) {winds and\\jets};
 \node[source,below right=of stars] (sn) {supernovae and\\cosmic rays};
 \node[phase,below=17mm of winds] (cloud) {cloud disruption\\and compression};
 \node[phase,left=of cloud] (hot) {hot cavities and\\superbubbles};
 \node[phase,right=of cloud] (turb) {turbulence, fields,\\phase exchange};
 \node[phase,below=of cloud] (out) {fountain or wind:\\mass, metals, energy};
 \draw[arrow] (stars)--(rad); \draw[arrow] (stars)--(winds); \draw[arrow] (stars)--(sn);
 \draw[arrow] (rad)--(cloud); \draw[arrow] (winds)--(hot); \draw[arrow] (sn)--(hot);
 \draw[arrow] (sn)--(turb); \draw[arrow] (hot)--(cloud); \draw[arrow] (turb)--(cloud);
 \draw[arrow] (hot)--(out); \draw[arrow] (cloud)--(out); \draw[arrow] (turb)--(out);
\end{tikzpicture}
\caption{Original feedback coupling map.  Arrows represent transfer, not a strict
time sequence.  Observed phases exchange mass and energy, so summing independently
derived phase rates can double count the same material.}
\label{fig:feedback-coupling-map}
\end{figure}

\section{Cosmic rays and magnetic coupling}

Supernova shocks and other accelerators populate nonthermal particles.  Cosmic rays
contribute pressure, ionize shielded gas, heat through wave damping, and can drive
outflows if their pressure gradient couples to plasma.  Transport may involve
advection, diffusion, and streaming, each anisotropic in an ordered magnetic field.

Gamma rays probe hadronic interactions and energetic particles; synchrotron emission
depends on electron spectra and magnetic field; ion chemistry constrains low-energy
ionization.  These observables sample different portions of the cosmic-ray spectrum.
A single ``cosmic-ray energy density'' inferred from radio equipartition is not a
direct measurement of the particles that dominate molecular ionization.

Cosmic-ray-driven winds can accelerate gas more gradually than a purely thermal wind
and may remain effective where radiative cooling removes thermal energy.  Quantitative
predictions depend on transport coefficients and wave damping, which are not universal
constants.  Line diagnostics of ionization state, temperature, and kinematics provide
important indirect constraints.

\section{How multiphase outflows are observed}

Different tracers select different phases and geometries:
\begin{table}[tbp]
\centering
\caption{Outflow tracers and principal selection effects.}
\label{tab:outflow-tracers}
\begin{tabularx}{\textwidth}{@{}
 >{\raggedright\arraybackslash}p{0.19\textwidth}
 >{\raggedright\arraybackslash}p{0.27\textwidth}
 >{\raggedright\arraybackslash}X@{}}
\toprule
Phase or component & Representative diagnostics & Main selection effects \\
\midrule
Hot volume-filling wind & X-ray continuum and highly ionized lines & Emission weights
dense interfaces; absorption requires a bright background; temperature and abundance
degenerate with emission measure \\
Warm ionized & H$\alpha$, [N\,II], [S\,II], [O\,III], UV ions & Extinction, shocks
versus photoionization, density-squared weighting, projection and beam smearing \\
Warm neutral & Na\,I D, Mg\,II, H\,I absorption and emission & Resonant scattering,
stellar contamination, ionization/depletion corrections, covering fraction \\
Molecular & CO, OH, HCN, HCO$^+$, warm H$_2$ & Conversion factor, excitation,
opacity, abundance, interferometric filtering, orientation \\
Dusty gas & Far-IR continuum, extinction, scattered light & Dust temperature,
opacity, dust-to-gas ratio, foreground versus outflow geometry \\
Cosmic rays & Synchrotron, gamma rays, ionization chemistry & Magnetic field,
transport, loss time, unrelated AGN contribution \\
\bottomrule
\end{tabularx}
\end{table}

Absorption detects only gas projected against a background source, yielding a
covering-weighted pencil beam.  Emission can map the flow but strongly weights dense
clumps.  Resonant emission may be scattered into locations far from photon production.
Agreement of velocity alone does not prove that phases are co-spatial.

\section{Mass, momentum, and energy of an outflow}

For an ionized phase whose recombination luminosity is
$L\propto n_en_pV$, mass scales as
\begin{equation}
 M_{\rm ion}\simeq
 \frac{\mu m_{\rm H}L_{\rm rec}}
 {h\nu\alpha^{\rm eff}(T)n_e},
 \label{eq:ionized-outflow-mass}
\end{equation}
after extinction and contamination corrections.  Density uncertainty maps directly
to mass uncertainty.  Density-sensitive doublets often measure bright clumps rather
than the low-density material that may dominate volume or mass.

For a thin partial shell at radius $R$ subtending solid angle $\Omega$, an absorption
column gives
\begin{equation}
 M_{\rm shell}\simeq\Omega R^2\mu m_{\rm H}N_{\rm H},
 \label{eq:shell-mass-column}
\end{equation}
where $N_{\rm H}$ requires ionization, abundance, depletion, and covering corrections.
The radius is often the least directly constrained quantity for unresolved absorption.

Several incompatible conventions are used for outflow rate.  For a shell crossing
time $R/v$,
\begin{equation}
 \dot M_{\rm out}\simeq\frac{M_{\rm out}v}{R}.
 \label{eq:outflow-rate-thin-shell}
\end{equation}
A uniformly filled steady cone may introduce a factor of order a few.  Rather than
calling one convention correct in every geometry, authors should provide $M$, $v$,
$R$, and the exact equation so readers can translate.

Velocity choices are likewise consequential.  A centroid measures bulk shift; a
maximum or percentile velocity emphasizes wings; deprojected speed depends on cone
inclination.  Because energy scales as $v^2$, wing definitions dominate kinetic-power
comparisons.

\subsection{Worked example: geometry dominates an outflow rate}

Suppose a molecular outflow has $M_{\rm out}=10^8M_\odot$, characteristic speed
$v=300\,\mathrm{km\,s^{-1}}$, and radius $R=1$~kpc.  The flow-time estimate is
\begin{equation}
 \dot M_{\rm out}\simeq31
 \left(\frac{M_{\rm out}}{10^8M_\odot}\right)
 \left(\frac{v}{300\,\mathrm{km\,s^{-1}}}\right)
 \left(\frac{R}{1\,\mathrm{kpc}}\right)^{-1}
 M_\odot\,\mathrm{yr^{-1}}.
\end{equation}
If the emission is actually concentrated at 0.3~kpc, the rate rises by more than a
factor of three.  If the molecular conversion factor shifts by a factor of four, so
does the rate.  A formal line-flux uncertainty of five percent is not the relevant
total uncertainty.

\section{Launch, survival, and phase creation}

Cool clouds placed in a fast hot wind experience ram pressure, shocks, shear, thermal
conduction, magnetic draping, and radiative cooling.  A characteristic cloud-crushing
time is
\begin{equation}
 t_{\rm cc}=\frac{\chi^{1/2}r_c}{v_{\rm rel}},
 \label{eq:cloud-crushing-time}
\end{equation}
where $\chi=\rho_c/\rho_h$.  In simple nonradiative hydrodynamics, disruption can occur
before acceleration to the hot-wind speed.  Radiative cooling in mixed gas, magnetic
fields, conduction, self-gravity, and continuous cloud creation can extend survival
or create new cool material in situ.

Therefore cool gas at large radius need not be the same parcels launched cold from
the disk.  It may have condensed from a mixed outflow.  Metallicity, dust survival,
ion ratios, cloud size distribution, and velocity alignment can test origin, but no
single diagnostic is decisive.

\section{Escape, fountains, and recycling}

Comparing line-of-sight velocity with an escape speed is only a first check.  The gas
has a three-dimensional velocity distribution, the gravitational potential evolves,
and hydrodynamic drag, cooling, and pressure forces continue after observation.  A
parcel below local escape speed can be accelerated; one above it can transfer energy
and fall back.

Material that returns in a galactic fountain still matters: it redistributes metals
and angular momentum, exchanges mass with the halo, and delays star formation.  The
CGM can contain first-time inflow, recycled wind, stripped satellite gas, and material
circulating within the halo \citep{TumlinsonPeeplesWerk2017}.  An absorption component
labeled ``outflow'' from its velocity sign alone may be ambiguous without galaxy
orientation and systemic-velocity precision.

\section{A rigorous feedback claim}

A strong feedback analysis proceeds through five tests:
\begin{enumerate}
 \item \textbf{Identification:} show that the spectral component is distinct from
rotation, tidal structure, beam smearing, and unrelated foreground gas.
 \item \textbf{Excitation:} determine whether photoionization, shocks, mixing, or an
AGN supplies the observed lines; allow mixtures.
 \item \textbf{Geometry:} state the aperture, radius, solid angle, inclination,
covering fraction, and whether absorption and emission sample the same structure.
 \item \textbf{Budget:} derive mass, momentum, and energy with conversion-factor,
density, ionization, extinction, and phase corrections made explicit.
 \item \textbf{Causality and fate:} compare source history with flow time and assess
whether gas is dispersed, recycled, or escaping rather than equating detection with
quenching.
\end{enumerate}

The ideal output is not one ``feedback efficiency'' with a tiny statistical error.
It is a phase-resolved budget with transparent conditional ranges.  Such a budget can
be compared meaningfully with simulations, which must apply the same apertures, line
emissivities, and selection functions to their synthetic observations.
